\documentclass[11pt]{article}

\usepackage[margin=1in]{geometry}
\usepackage{amsmath,amssymb,amsthm,mathtools}
\usepackage{enumitem}
\usepackage{booktabs}
\usepackage{microtype}
\usepackage[T1]{fontenc}

\newtheorem{conjecture}{Conjecture}
\newtheorem{lemma}{Lemma}
\newtheorem{proposition}{Proposition}
\theoremstyle{definition}
\newtheorem{definition}{Definition}
\theoremstyle{remark}
\newtheorem{remark}{Remark}

\newcommand{\wt}{\operatorname{wt}}
\newcommand{\Peak}{\operatorname{Peak}}
\newcommand{\PPeak}{\operatorname{PPeak}}
\newcommand{\Dbl}{\operatorname{Dbl}}
\newcommand{\Diag}{\operatorname{Diag}}
\newcommand{\RR}{\operatorname{RR}}
\newcommand{\cB}{\mathcal{B}}
\newcommand{\cQ}{\mathcal{Q}}
\newcommand{\cR}{\mathcal{R}}
\newcommand{\cP}{\mathcal{P}}

\title{The Shifted $Q/R$ and $P$ Conjectures and Their Exhaustive Small-Case Checks}
\author{}
\date{}

\begin{document}
\maketitle

\section{Shifted set-valued tableaux}

A \emph{strict partition} is a finite sequence
\[
  \lambda=(\lambda_1>\lambda_2>\cdots>\lambda_{\ell}>0).
\]
Its shifted diagram is
\[
  D_\lambda=\{(r,c):1\le r\le \ell,\ r\le c\le r+\lambda_r-1\}.
\]
If $\mu\subseteq\lambda$ are strict partitions, then
$D_{\lambda/\mu}=D_\lambda\setminus D_\mu$ is a skew shifted diagram.  A box
$(r,r)$ is called diagonal.

We use the marked alphabet
\[
  1'<1<2'<2<3'<3<\cdots.
\]
A shifted set-valued tableau assigns a nonempty finite set of letters to each
box.  It is semistandard when the boxes are weakly increasing from left to
right and from top to bottom in the set-valued sense, no primed letter $i'$ is
repeated in one row, and no unprimed letter $i$ is repeated in one column.  Its
weight is
\[
  \wt(T)=(m_1,m_2,\ldots),
\]
where $m_i$ is the total number of occurrences of $i$ and $i'$.

There are three classes of tableaux.
\begin{itemize}[leftmargin=2em]
\item A $Q$-tableau is a tableau satisfying the conditions above, with no
additional diagonal restriction.
\item An $R$-tableau is a $Q$-tableau satisfying the $i$-primed property for
every $i$: in the first box in row-reading order containing $i$ or $i'$, the
letter $i'$ occurs and the letter $i$ does not occur.
\item A $P$-tableau is a $Q$-tableau in which every entry in a diagonal box is
primed.
\end{itemize}
The $P$ convention is discussed again in Section~\ref{sec:diagonal-choice}.

The row-reading order $\RR$ reads rows from bottom to top and boxes from left
to right within a row.  For the lattice test, the entries inside one box are
read in the special order
\[
  1<2<3<\cdots<\cdots<3'<2'<1';
\]
that is, unprimed entries in increasing order followed by primed entries in
decreasing order.  This special order is used only to form the lattice reading
word.

Fix $i$.  The $i$-lattice test compares the letters $i,i',(i+1),(i+1)'$.
Start with equal tallies for $i/i'$ and $(i+1)/(i+1)'$.  First scan the reading
word backward.  Count the unprimed letters and never allow the $(i+1)$ tally
to exceed the $i$ tally; moreover, $(i+1)'$ is forbidden when the two tallies
are tied.  Retaining these tallies, scan the word forward.  Count the primed
letters and again never allow the upper tally to exceed the lower tally;
moreover, an unprimed $i$ is forbidden at a tie.  A tableau is \emph{lattice}
when this holds for every $i$.

Write $GQ_{\lambda/\mu}$, $GR_{\lambda/\mu}$, and $GP_{\lambda/\mu}$ for the
weight generating functions of the corresponding tableaux:
\[
  G\!X_{\lambda/\mu}=\sum_T x^{\wt(T)},
  \qquad X\in\{Q,R,P\}.
\]
The same notation without an inner shape denotes a straight shifted shape.

\section{The two conjectures}

For a skew shape $\lambda/\mu$ and a strict partition $\nu$, define
\[
 b_{\lambda/\mu}^{\nu}
 =\#\{T\in\cR(\lambda/\mu):T\text{ is lattice and }\wt(T)=\nu\},
\]
and
\[
 c_{\lambda/\mu}^{\nu}
 =\#\{T\in\cP(\lambda/\mu):T\text{ is lattice and }\wt(T)=\nu\}.
\]
Thus the $R$ coefficient uses both the $i$-primed property and the lattice
property, while the $P$ coefficient uses the diagonal condition and the same
lattice property.

\begin{conjecture}[$Q/R$ conjecture]
For every skew shifted shape $\lambda/\mu$,
\[
  GQ_{\lambda/\mu}
  =\sum_{\nu} b_{\lambda/\mu}^{\nu}GQ_\nu.
\]
Equivalently, after the reduction in Section~\ref{sec:budding-inflation}, the
same coefficients satisfy the corresponding $R$ identity.
\end{conjecture}

\begin{conjecture}[$P$ conjecture]
For every skew shifted shape $\lambda/\mu$,
\[
  GP_{\lambda/\mu}
  =\sum_{\nu} c_{\lambda/\mu}^{\nu}GP_\nu.
\]
\end{conjecture}

The coefficient weights are automatically strict.

\begin{lemma}[Strict dominance]
A nonzero lattice $R$- or $P$-tableau has strictly dominant weight.
\end{lemma}

\begin{proof}
Suppose that the total numbers of $i/i'$ and $(i+1)/(i+1)'$ are equal and
nonzero.  Let $B$ be the last box in $\RR$ order containing any of
\[
  i,\quad i+1,\quad (i+1)',\quad i',
\]
and let $S$ be the intersection of this set with $B$.  In the special order
inside $B$, the relevant forward order is
\[
  i,\ i+1,\ (i+1)',\ i'.
\]
If $i'\in S$, then immediately before the final $i'$ is read in the forward
scan, the upper tally is ahead.  If $(i+1)'\in S$, it is the first relevant
letter in the backward scan and is read at a tie.  After excluding these two
letters, an $i+1$ in $S$ would be first in the backward scan and would put the
upper tally ahead.  Finally, if $i\in S$, it is the last relevant letter in the
forward scan, when all relevant letters have been accounted for and the
tallies are tied.  Every possibility is forbidden, so $S$ is empty, contrary
to the choice of $B$.
\end{proof}

\section{Budding and inflation in the $Q/R$ case}
\label{sec:budding-inflation}

\subsection{Budding}

For every value $i$ appearing in an $R$-tableau, the first $i'$ may be changed
independently into
\[
  i',\qquad i,\qquad \{i',i\}.
\]
The first two choices preserve the exponent of $x_i$, while the third raises
it by one.  Hence define the linear budding operator on monomials by
\[
  \cB(x^\alpha)
  =\prod_{\alpha_i>0}\left(2x_i^{\alpha_i}+x_i^{\alpha_i+1}\right).
\]
For every straight or skew shifted shape $\theta$,
\[
  GQ_\theta=\cB(GR_\theta).
\]
Consequently, an $R$ expansion with coefficients
$b_{\lambda/\mu}^{\nu}$ buds to the $Q$ expansion with exactly the same
coefficients.  There is therefore no separate finite ``$GR$ check'' and
``$GQ$ check'' in the program: both are represented by the same final
combinatorial equality below.

\subsection{Inflation and the peak-refined equality}

A standard shifted set-valued tableau of size $n$ uses each label
$1,2,\ldots,n$ exactly once.  In the inflated class used for the $Q/R$ check,
consecutive labels are not allowed to occupy the same box.  Inflation replaces
a semistandard tableau by such a standard object; conversely, the possible
semistandardizations of an inflated object depend only on its peak set.

For a standard tableau $S$, let $b_j$ be the box containing $j$.  Define
\[
  \Peak(S)=\{j\in\{2,\ldots,n-1\}:
  b_{j-1}\text{ is strictly left of }b_j
  \text{ and }b_{j+1}\text{ is strictly below }b_j\}.
\]
For a shape $\theta$, an integer $n$, and a peak set $P$, let
\[
  a_\theta(n,P)
  =\#\{S:\operatorname{shape}(S)=\theta,\ |S|=n,
           \Peak(S)=P,
           \text{ no consecutive labels share a box}\}.
\]
The final combinatorial equality on which the $Q$ conjecture relies is
\begin{equation}
\label{eq:qr-final}
  a_{\lambda/\mu}(n,P)
  =\sum_\nu b_{\lambda/\mu}^{\nu}a_\nu(n,P)
\end{equation}
for every $n$ and every peak set $P$.  Budding and inflation reduce the
original $Q$ identity to exactly these peak-refined equalities.

There is a parallel reduction for the $P$ conjecture, but it requires more
statistics than an ordinary peak set and it compares dominant monomials rather
than the raw refined records.  The precise rule used by the code is given in
Section~\ref{sec:p-reduction}.

\section{What the code generates}

The file \texttt{code.py} contains both checks.  The bound is configurable;
\texttt{--mode qr}, \texttt{--mode p}, and \texttt{--mode both} select the
check.

\subsection{Inductive generation}

The program builds standard set-valued shifted tableaux by inserting the next
largest label.  From a tableau of shape $\rho$, the new label is placed either
in a new addable outer-corner box or in an existing terminal box.  A terminal
box is the rightmost box of a row with no box immediately below it.  These are
exactly the positions in which the new largest label can be inserted while
preserving the shifted semistandard conditions.

For the $Q/R$ check, insertion in the same box as the preceding label is
rejected.  For the $P$ check, it is retained.

Each compressed state stores only the information needed for future growth:
the current shape, the boxes containing the last two labels, the relevant
statistic masks, and the canonical-prefix data.  States with the same stored
data are merged and carried with a multiplicity.

The canonical standard tableau on a shifted shape has one entry per box and is
filled by consecutive labels along rows, from left to right and from the top
row downward.  For a generated tableau $S$, the program tracks the largest
initial canonical region, together with all its canonical initial subshapes.
If a later label is inserted into one of its singleton boxes, the canonical
chain is truncated at that point.

Whenever a new straight tableau with $N$ labels is produced, every nonempty
canonical initial subshape $\mu$ of size $x$ gives a skew tableau by deleting
the labels $1,\ldots,x$.  The remaining labels are shifted down by $x$.  Thus
\[
  x+(N-x)=N,
\]
so one inductive generation simultaneously produces every straight object of
size at most the bound and every skew object for which
\[
  \#\text{missing boxes}+\#\text{actual entries}
  \le \text{bound}.
\]
A peak index is shifted down by $x$; if it becomes $1$, it is discarded because
the smallest remaining label cannot be a peak.  The other index sets in the
$P$ check are shifted in the same way, with index $0$ discarded.

\subsection{The coefficient side}

For each skew shape, the program independently enumerates the flagged
semistandard set-valued tableaux needed for the conjectured coefficients.  A
letter $i'$ is encoded by $2i-1$ and $i$ by $2i$.  It checks the row and column
conditions, forms the $\RR$ word using the special within-box order, and then
applies the lattice test.

For the $Q/R$ coefficients it also applies the $i$-primed property.  For the
$P$ coefficients it requires every entry in a present diagonal box to be
primed.  The surviving tableaux are counted by weight.  The strict-dominance
lemma is also checked internally: a surviving nonempty weight that is not a
strict partition raises an error.

\subsection{The $Q/R$ comparison}

For every generated skew shape $\lambda/\mu$, every permitted actual-entry
number $n$, and every peak set $P$, the code compares the left side of
\eqref{eq:qr-final} with the coefficient-weighted straight side.  Since the
inflation reduction has already replaced semistandard tableaux by the
no-consecutive standard objects, no additional monomial transformation is
needed in this part of the program.

\section{The $P$ reduction used in the code}
\label{sec:p-reduction}

The standard tableaux used for the $P$ check allow consecutive labels in one
box.  Let $b_j$ be the box containing $j$.  For distinct boxes, write
$b_1\nearrow b_2$ if $b_2$ is northeast of $b_1$, and write
$b_1\swarrow b_2$ if $b_2$ is southwest of $b_1$; equality means that the two
labels are in the same box.

For three consecutive labels $j-1,j,j+1$, the middle label $j$ is a
$P$-peak when
\[
  b_{j-1}\ (\nearrow\text{ or }=)\ b_j
  \qquad\text{and}\qquad
  b_j\ (\swarrow\text{ or }=)\ b_{j+1}.
\]
Equivalently, the four allowed relation pairs are
\[
  (\nearrow,\swarrow),\qquad (=,=),\qquad (=,\swarrow),
  \qquad(\nearrow,=).
\]
Define three statistics:
\begin{align*}
  \PPeak(S)&=\{j:j\text{ is a $P$-peak}\},\\
  \Dbl(S)&=\{j:b_j=b_{j+1}\},\\
  \Diag(S)&=\{j:b_j\text{ is diagonal}\}.
\end{align*}
Thus $\Dbl(S)$ records every same-box consecutive pair, not only those in a
diagonal box.

The code first counts standard tableaux by
\[
  (n,\PPeak(S),\Dbl(S),\Diag(S)).
\]
It then converts this counter into a counter of dominant monomials.  Let
$\alpha=(\alpha_1,\ldots,\alpha_r)\vdash n$.  Interpret its parts as consecutive
intervals
\[
 I_1=[1,\alpha_1],\quad
 I_2=[\alpha_1+1,\alpha_1+\alpha_2],\quad\ldots.
\]
The partition $\alpha$ \emph{fits} the record of $S$ when
\begin{enumerate}[label=(\roman*),leftmargin=2.5em]
\item every $P$-peak is the first or last label of one of the intervals $I_k$;
\item every $d\in\Dbl(S)\cap\Diag(S)$ is the last label of an interval.
\end{enumerate}
For a fitting partition, call $I_k$ free when it contains no label in
$\Diag(S)$ and does not contain both $d$ and $d+1$ for any
$d\in\Dbl(S)$.  If $f(S,\alpha)$ is the number of free parts, then $S$
contributes the dominant monomial indexed by $\alpha$ with multiplicity
\[
  2^{f(S,\alpha)}.
\]
The factor records the independent prime/unprime choice for each part not
forced by a diagonal occurrence or by a same-box consecutive pair contained
inside that part.

For each skew shape and degree, the program applies this transformation to the
skew records and to the coefficient-weighted straight records and compares the
resulting dominant-monomial counters.  The functions under consideration are
symmetric, so equality of these dominant coefficients is the required finite
check.  Raw equality of the quadruple-refined records is neither asserted nor
needed.

\section{The bound-$12$ computation}

The combined program was run with
\[
  \texttt{python code.py --mode qr --max-entries 12}
\]
and
\[
  \texttt{python code.py --mode p --max-entries 12}.
\]
Both checks passed.  The recorded totals were
\begin{center}
\begin{tabular}{lrr}
\toprule
 & $Q/R$ & $P$\\
\midrule
straight shapes & 69 & 69\\
proper skew shapes & 1172 & 1172\\
straight standard tableaux & 42050 & 583251\\
derived skew standard tableaux & 146283 & 1116632\\
shape--degree comparisons & 7213 & 7213\\
\bottomrule
\end{tabular}
\end{center}

More explicitly, for every proper skew shifted shape $\lambda/\mu$ satisfying
\[
  |\mu|+|\lambda/\mu|\le 12,
\]
the code checks every actual-entry degree $n$ with
\[
  |\lambda/\mu|\le n\le 12-|\mu|.
\]
In the $Q/R$ case, it checks every peak-refined equality
\eqref{eq:qr-final}; by budding and inflation, this verifies the conjectured
$Q$ expansion throughout that finite range.  In the $P$ case, it checks every
dominant monomial obtained from the fitting rule above; this verifies the
conjectured $P$ expansion throughout the same finite range.  These are
exhaustive checks under the charged bound, not proofs for arbitrary shapes or
degrees.

\section{Why diagonal entries are taken to be primed}
\label{sec:diagonal-choice}

One may equivalently formulate $P$-tableaux with no primed entries on the
diagonal.  Replacing every diagonal $i$ by $i'$ preserves the lattice
property.  For a fixed pair $i,i+1$, the only relevant diagonal patterns are:
only $i$, only $i+1$, both in one diagonal box, or $i$ in a diagonal box and
$i+1$ in the next southeast diagonal box.  The first three cases are immediate
from the two tally scans.  In the last case, one considers the portion of the
two rows extending through the last $i+1$.  The apparent one-unit difference
in the backward scan is exactly the forbidden forward tie at the diagonal
$i$, so the two conventions again agree.

We use the all-primed diagonal convention because it states the $P$ condition
as a restriction on unprimed entries.  The $R$ condition is likewise naturally
phrased by excluding an unprimed $i$ from the first relevant box.  The chosen
convention therefore keeps the priming restrictions in the $Q/R$ and $P$
settings parallel without changing either the lattice tableaux or the
conjectured coefficients.

\end{document}
